% ==============================================================================
% PAGE 42: CHAPTER-7 WORKING (Numbered Page 34)
% ==============================================================================
\begin{center}
  {\fontsize{15}{18}\selectfont \bfseries \textcolor{red!80!black}{\underline{CHAPTER-7}}}\\[2mm]
  {\fontsize{16}{20}\selectfont \bfseries \textcolor{red!80!black}{\underline{WORKING}}}\\[5mm]
\end{center}

\noindent
\begin{spacing}{1.2}
{\fontsize{10}{14}\selectfont
The working of the Parabolic Solar Water Heater with Automatic Tracking and Power Generation is based on the principle of solar energy concentration and automatic solar tracking. The parabolic dish collects and focuses sunlight onto a single focal point. The surface of the dish is covered with mirror pieces, which act as reflective material to direct the sunlight accurately toward the center. A black-painted copper receiver tube is fixed at this focal point to absorb the maximum heat. When water passes through the receiver tube, it gains thermal energy from the concentrated solar rays, causing a steady rise in its temperature. The setup ensures that maximum solar radiation is utilized for heating the water efficiently.\\[3.5mm]
The automatic solar tracking mechanism allows the parabolic dish to continuously face the sun from morning to evening. Two Light Dependent Resistors (LDRs) are mounted on opposite sides of the dish and connected to the control circuit. These sensors detect the difference in sunlight intensity between the two sides. When one sensor receives more light, the control circuit sends a signal to the microcontroller and with help of motor driver motor is start rotating, which rotates the dish slightly---about 4 to 5 degrees at a time---until both LDRs receive equal sunlight. This gradual movement keeps the dish properly aligned with the sun throughout the day, maintaining the correct focus on the receiver and improving overall heat absorption efficiency.\\[3.5mm]
The electrical power system of the project operates through a 6V solar-powered battery. A solar panel converts sunlight into electrical energy and charges the 6V battery via a charge controller. The stored electrical energy is used to power the PMDC motor, LDR sensors, and the motor driver circuit for automatic tracking. A blocking diode is used in the circuit to allow one-way current flow, preventing reverse discharge of the battery at night. The combined operation of the solar panel, battery, sensors, and motor ensures that the system runs automatically without manual adjustment. Thus, the parabolic solar heater continuously tracks the sun, heats the water efficiently, and uses clean solar energy both for heating and for controlling its own movement.
}
\end{spacing}
\clearpage
